\documentclass{article}
\input{../../lib.sty}

\title{\texttt{fn choose}}
\author{Michael Shoemate}
\date{}

\begin{document}
\maketitle

Proves soundness of \texttt{fn choose} in \texttt{mod.rs}.

\section{Hoare Triple}
\subsection*{Precondition}
None.

\subsection*{Pseudocode}
\lstinputlisting[language=Python,firstline=2,escapechar=|]{./pseudocode/choose.py}

\subsection*{Postcondition}
\begin{theorem}
\label{postcondition}
For any \texttt{best} and \texttt{next}, \texttt{choose(best, next)} returns:
\begin{itemize}
    \item \texttt{best} unchanged if \texttt{next[0]} is NaN;
    \item \texttt{next} if \texttt{best = None} and \texttt{next[0]} is not NaN;
    \item otherwise the higher-scoring of \texttt{best} and \texttt{next}, breaking ties in favor
    of \texttt{next}.
\end{itemize}
\end{theorem}

\section{Proof}
The function has three branches:
\begin{itemize}
    \item line \ref{line:nan-guard} returns \texttt{best} unchanged on NaN;
    \item line \ref{line:best-none} returns \texttt{next} when there is no current best;
    \item line \ref{line:compare} returns \texttt{best} only when its score is strictly larger,
    and otherwise the final return chooses \texttt{next}, which implements the stated tie-break.
\end{itemize}

\end{document}
